\section{Introduction}
\label{sec:introduction}

% Goals of this section (course writing recs):
%  - state the research question
%  - why it is important / interesting
%  - limitations of prior approaches (full-item translation, question-side code-switch)
%  - our approach (answer-side-only translation)
%  - benefits + the main experiments
%
% RECOMMENDED: a figure/example here. Suggestion: one CSQA item rendered across the
% four conditions (EN-EN / EN-X / X-X / X-EN) with the gold answer highlighted, so a
% reader instantly sees what "answer-side translation" means and what a "flip" is.

% Research question: Do language models choose the correct commonsense \emph{concept},
% or do they rely on English answer \emph{wording}?

% TODO: prose.

% \begin{figure}[t]
%   \centering
%   \includegraphics[width=\columnwidth]{figures/teaser.pdf}
%   \caption{One CommonsenseQA item across the four answer-language conditions.
%            The question stays English; only the choices change language. A
%            \emph{flip} is a change in the selected choice across conditions.}
%   \label{fig:teaser}
% \end{figure}
